% EM structure and teaching revision 20260928
\subsection{混合频率信号的有效值与适用范围}
设周期信号含直流和不同整数倍频率的谐波：
\[
 u(t)=U_d+\sum_{n=1}^{N}U_{mn}\cos(n\omega t+\phi_n).
\]
在共同周期内，直流与各正弦项的乘积平均为零，不同频率项之间的乘积平均也为零，因此
\[
 U_{\mathrm{rms}}^2=U_d^2+\frac12\sum_{n=1}^{N}U_{mn}^2.
\]
同频项却不能按平方直接相加，必须先考虑相位。两项幅值同为 $A$ 的同频电压，相位相同会合成峰值 $2A$，反相则完全抵消。
例如 $u=3+4\cos\omega t+2\sin2\omega t$（单位为伏），平均值为 $3\,\mathrm V$，有效值为 $\sqrt{19}\,\mathrm V$。若接在线性电阻 $R$ 上，平均发热功率为 $19/R$；若接在频率相关网络上，则须分别求各谐波支路响应后再合成实际波形。
相量法一次只处理一个频率。把多频信号压缩成一个“总相量”，会丢掉各频率面对的不同感抗、容抗及传播条件。

